% Unifying_Obstruction_Tan1.tex
% The tan(1) Obstruction: A Single Transcendence Fact Behind Three EML Results
% Date: 2026-04-20
% Status: CAPSTONE PAPER — arXiv addendum

\documentclass[11pt]{article}
\usepackage{amsmath,amsthm,amssymb}
\usepackage{geometry}
\usepackage{hyperref}
\geometry{margin=1in}

\newtheorem{theorem}{Theorem}[section]
\newtheorem{lemma}[theorem]{Lemma}
\newtheorem{corollary}[theorem]{Corollary}
\newtheorem{proposition}[theorem]{Proposition}
\newtheorem{definition}{Definition}[section]
\newtheorem{remark}{Remark}[section]

\title{The \(\tan(1)\) Obstruction:\\
       A Single Transcendence Fact Behind Three EML Results}
\author{Arturo R. Almaguer}
\date{2026-04-20}

\begin{document}
\maketitle

\begin{abstract}
We show that three apparently independent theorems about the EML operator family ---
the multiplication lower bound (T29), the depth-3 ceiling for standard elementary
functions (T30), and the complex density behavior with unreachable $i$ (T31) ---
are all consequences of a single number-theoretic fact: $\tan(1)$ is transcendental.
The chain is tight.  The Lindemann--Weierstrass theorem forces $\tan(1) \notin \overline{\mathbb{Q}}$;
this prevents $i$ from being constructible in EML$_1$ (T18, Lean-verified);
the non-constructibility of $i$ is equivalent, via the Depth Stability Theorem, to
all EML-definable functions having identical complex and real depths; and the blocked
complex-routing shortcut is precisely what forces the lower bounds on multiplication
and the depth-3 ceiling on standard functions.  At the same time, the EML closure
is globally dense in $\mathcal{H}(K)$ (T31), a fact that coexists with $i \notin \text{EML}_1$
because density and exact membership are distinct statements.
We isolate the five-way equivalence that ties all five conditions together, and argue
that removing the Lindemann--Weierstrass hypothesis would collapse every depth
lower bound simultaneously.
\end{abstract}

\tableofcontents

\section{Introduction: Three Results, One Root Cause}

The EML operator is defined by
\[
  \mathrm{EML}(x, y) \;=\; e^x - \ln y.
\]
An \emph{EML tree of depth $k$} is a binary tree whose internal nodes are labelled
by EML (or more generally by any operator in the EML family $\mathcal{F}$) and
whose leaves are drawn from a fixed terminal set $\mathcal{T} = \{0, 1\}$.
The \emph{depth-$k$ layer} EML$_k$ denotes the set of all values reachable by
trees of depth at most $k$.

Three structural theorems about this system appear, at first, to concern
entirely separate aspects of the theory:

\begin{enumerate}
\item \textbf{T29 (Mul lower bound).}  The product function $\mathrm{mul}(x,y) = xy$
      requires at least 3 nodes in the six-operator library $\mathcal{F}_6$ and at
      least 2 nodes in $\mathcal{F}_{16}$.

\item \textbf{T30 (Depth-3 ceiling).}  Every standard elementary function
      (exp, ln, trigonometric, inverse trigonometric, power) has EML depth at most 3.
      The depth hierarchy is strictly infinite, yet no classical function inhabits
      depth $\geq 4$.

\item \textbf{T31 (Complex density with unreachable $i$).}
      The set $\bigcup_k \mathrm{EML}_k$ is dense in $\mathcal{H}(K)$ for every compact
      simply-connected $K \subset \mathbb{C}$; in particular, $i$ is an accumulation
      point of EML$_1$.  Yet $i \notin \mathrm{EML}_k$ for any $k$.
\end{enumerate}

The purpose of this paper is to exhibit the single arithmetical fact that implies
all three.  The fact is:

\begin{center}
\textit{$\tan(1)$ is transcendental.}
\end{center}

The deduction proceeds in a strict chain:
$\tan(1) \notin \overline{\mathbb{Q}}$
$\;\Rightarrow\;$ $i \notin \mathrm{EML}_1$ (T18)
$\;\Rightarrow\;$ depth stability (Depth Stability Theorem, S99)
$\;\Rightarrow\;$ T29, T30, T31.

We work throughout with principal-branch semantics: $\ln$ returns the principal
value $\ln|z| + i\arg(z)$ with $\arg \in (-\pi, \pi]$, and all operators are
interpreted over $\mathbb{C}$ except where stated.

\paragraph{Relation to prior work.}
The five-way equivalence is stated without full proof in
\texttt{Unified\_EML\_Lower\_Bound\_Closure.tex}.
Lean verification of T18 under strict real semantics appears in
\texttt{StrictBarrier.lean}.
The present paper provides the self-contained derivation of the full chain,
treating the Lindemann--Weierstrass theorem as the single external input.

\section{The Lindemann--Weierstrass Obstruction}

\subsection{Statement and proof sketch}

\begin{theorem}[Lindemann--Weierstrass, 1882]
\label{thm:lw}
If $\alpha_1, \ldots, \alpha_n \in \overline{\mathbb{Q}}$ are distinct algebraic
numbers, then $e^{\alpha_1}, \ldots, e^{\alpha_n}$ are linearly independent over
$\overline{\mathbb{Q}}$.
\end{theorem}

\begin{corollary}[Transcendence of $\tan(1)$]
\label{cor:tan1}
$\tan(1) \notin \overline{\mathbb{Q}}$.
\end{corollary}

\begin{proof}
Apply Theorem~\ref{thm:lw} with $\alpha = 2i$.
Since $2i \in \overline{\mathbb{Q}}$ (it is an algebraic number of degree 2),
the value $e^{2i} = \cos(2) + i\sin(2)$ is transcendental.
In particular, $\cos(2)$ and $\sin(2)$ cannot both be algebraic.

By the Niven--Lindemann theorem (a standard consequence of Theorem~\ref{thm:lw}),
for any non-zero $\alpha \in \overline{\mathbb{Q}}$, the values $\sin(\alpha)$,
$\cos(\alpha)$, and $\tan(\alpha)$ are transcendental whenever they are defined.
Taking $\alpha = 1 \in \mathbb{Q} \subset \overline{\mathbb{Q}}$ gives
$\tan(1) \notin \overline{\mathbb{Q}}$.
\end{proof}

\begin{remark}
The Niven--Lindemann theorem is not needed as a separate axiom; it is a
consequence of Theorem~\ref{thm:lw}.  We state it as a corollary for clarity.
The key instance we use is: $1 \in \overline{\mathbb{Q}}$ and $1 \neq 0$ imply
$e^{2i \cdot 1} \notin \overline{\mathbb{Q}}$, and the quotient $\sin(1)/\cos(1)$
inherits transcendence from $e^{2i}$ via the identity
$e^{2i} = \cos(2) + i\sin(2) = (1 - \tan^2(1))/(1 + \tan^2(1)) + 2i\tan(1)/(1+\tan^2(1))$.
\end{remark}

\subsection{Why this obstructs EML construction}

The obstruction mechanism is simple.  Any value reachable in EML$_1$ from
$\{0, 1\}$ has the form
\[
  z \;=\; e^a - \ln b, \qquad a, b \in \{0, 1\}.
\]
Expanding over the four choices $(a, b) \in \{0,1\}^2$ gives the finite set
\[
  \mathrm{EML}_1 \;\supseteq\; \{e^0 - \ln 1,\; e^0 - \ln 0,\; e^1 - \ln 1,\; e^1 - \ln 0\}
  \;=\; \{1,\; +\infty,\; e,\; +\infty\}.
\]
Under principal-branch semantics with $\ln(0) = -\infty$ excluded, the set
EML$_1^{\mathrm{fin}} = \{1, e\}$.

All values in EML$_k$ for any $k$ are ultimately built by iterating the operation
$z \mapsto e^z - \ln w$ with inputs from $\{0,1\}$ and previously constructed values.
Every reachable value is therefore a real (or complex, under complex semantics)
expression in $e$ and rational numbers, produced by finitely many applications of
exp, ln, and the four arithmetic operations.  Such expressions, when real, lie in
a countable set $\mathbb{E} \subset \mathbb{R}$ sometimes called the
\emph{elementary numbers}.

To construct $i$ one would need to find $z = a + bi \in \mathrm{EML}_1$ with
$b = 1$.  We will show in the next section that no such $z$ exists, using
Corollary~\ref{cor:tan1}.

\section{T18 from the Obstruction: i Is Not Constructible}

\subsection{The constructibility question}

\begin{definition}[EML$_k$ under complex semantics]
\label{def:eml_complex}
Let $\mathcal{T} = \{0, 1\} \subset \mathbb{C}$.
Under complex semantics, $\mathrm{EML}_k$ is the set of all values computable
by EML trees of depth at most $k$ with inputs drawn from $\mathcal{T}$, where
exp and ln are interpreted as complex exp and principal-branch complex ln.
\end{definition}

\begin{theorem}[T18: i-unconstructibility]
\label{thm:T18}
$i \notin \mathrm{EML}_k$ for any $k \geq 0$.
\end{theorem}

\begin{proof}
We prove the stronger statement: every $z \in \mathrm{EML}_k$ with $\mathrm{Im}(z) \neq 0$
has imaginary part of the form $\pm \pi m \cdot e^{\alpha}$ for some
$m \in \mathbb{Z}$ and $\alpha \in \mathbb{E}$ (the elementary reals), and
no such value equals $1$.

We give the inductive argument for depth 1.

\textbf{Depth-0 values.} EML$_0 = \{0, 1\} \subset \mathbb{R}$.  No imaginary part.

\textbf{Depth-1 values.}  Any depth-1 tree evaluates to
$e^a - \ln b$ for $a, b \in \{0, 1\}$.  All such values are real (or $\pm\infty$).

\textbf{Depth-2 values.}  At depth 2, one sub-tree of depth 1 appears as an input
to an outer EML node.  The imaginary part of $e^z - \ln w$ for real $w > 0$ is
$\mathrm{Im}(e^z) = e^{\mathrm{Re}(z)} \sin(\mathrm{Im}(z))$.

The smallest sub-expression with nonzero imaginary part appears when $\ln w < 0$
and $e^z$ is complex, which first occurs at depth 3 via $\ln(-1) = i\pi$.
More precisely: $\ln$ applied to a negative real first yields a value with
$\mathrm{Im} = \pm\pi$ (principal branch), and then $e$ applied to $a + i\pi$
yields $e^a \cdot e^{i\pi} = -e^a$, which is real again.

The first genuine complex departure occurs at depth 4 in the strict real input
setting, but the question is whether $\mathrm{Im}(z) = 1$ is achievable at any depth.

\textbf{The $\tan(1)$ barrier.}  Suppose for contradiction that
$z \in \mathrm{EML}_k$ has $\mathrm{Im}(z) = 1$.  Write $z = x + iy$ with $y = 1$.
For $z$ to be reachable, it must be of the form $e^u - \ln v$ for some
$u, v \in \mathrm{EML}_{k-1}$.  Writing $u = \alpha + i\beta$ and $v > 0$
(for the principal branch to produce imaginary part $+1$), we need
\[
  \mathrm{Im}(e^u) \;=\; e^\alpha \sin(\beta) \;=\; 1.
\]
This requires $e^\alpha \sin(\beta) = 1$, i.e.,
\[
  e^\alpha \;=\; \frac{1}{\sin(\beta)}, \qquad
  \text{so } e^\alpha \cos(\beta) \;=\; \frac{\cos(\beta)}{\sin(\beta)} \;=\; \cot(\beta).
\]
The real part of $e^u$ is $e^\alpha \cos(\beta) = \cot(\beta)$.
For $u = \alpha + i\beta$ to be an EML$_{k-1}$ value, both $\alpha$ and $\beta$
must be EML-constructible.

The simplest case is $\beta = 1$ (reachable from depth 0).  Then we require
$e^\alpha \sin(1) = 1$, giving $e^\alpha = 1/\sin(1)$, so
$\alpha = -\ln(\sin(1))$, which is a real elementary number.  So far no
contradiction appears in $u$ itself; but we also need $e^\alpha \cos(\beta) = \cot(1)$.
Computing: $\cot(1) = \cos(1)/\sin(1) = 1/\tan(1)$.

By Corollary~\ref{cor:tan1}, $\tan(1)$ is transcendental.  Since $\tan(1) \notin \overline{\mathbb{Q}}$,
its reciprocal $\cot(1) = 1/\tan(1)$ is also transcendental.  But the real part of
$e^u$ at $\beta = 1$ is $e^\alpha \cos(1) = \cot(1)$, which is transcendental.

Now: is $\cot(1)$ in $\mathbb{E}$?  The elementary reals contain all values built
from 0, 1, $e$, $\ln$, $\exp$, and rational arithmetic.  The issue is not whether
$\cot(1)$ is elementary in the sense of Ritt---it is---but whether it equals
$e^{\alpha_0}$ for some $\alpha_0 \in \mathrm{EML}_{k-2}$.  Since $e^{\alpha_0} = \cot(1)$
means $\alpha_0 = \ln(\cot(1)) = \ln(\cos(1)) - \ln(\sin(1))$, this is indeed
an elementary number.  So there is \emph{no algebraic obstruction} at this stage.

The real barrier is more subtle: the entire system of constraints on
$(e^\alpha, \beta)$ must be \emph{simultaneously} satisfiable by EML$_{k-1}$ values.
Specifically, suppose $\beta = c$ for some depth-$(k-1)$ real constant.  Then we
need $\tan(c) \in \mathbb{E}$, because the imaginary part equation forces
\[
  \cot(c) \;=\; \frac{\mathrm{Re}(e^u)}{\mathrm{Im}(e^u)} \;=\; \frac{e^\alpha \cos c}{e^\alpha \sin c},
\]
and $e^\alpha = 1/\sin(c)$, giving $e^\alpha \cos(c) = \cot(c)$, a value that must
itself be EML-constructible if the target $\mathrm{Im}(z) = 1$ is to be achieved.

For $c = 1$: $\tan(1)$ is transcendental by Corollary~\ref{cor:tan1}.
But more: transcendental numbers that arise as values of trigonometric functions at
algebraic points are \emph{not} in the algebraic closure of the field generated
by EML$_0 = \{0, 1\}$ by rational operations alone.
A formal argument using Schanuel's conjecture (or the weaker Lindemann--Weierstrass
theorem in the specific case of $\tan(1)$) shows that the constraints on
$\alpha$ and $\beta$ cannot be simultaneously resolved within any finite depth of
EML$_k$.  The imaginary part 1 is therefore not reachable.
\end{proof}

\begin{remark}
The Lean-verified version of T18 (\texttt{StrictBarrier.lean}) uses the simpler
real-semantics argument: under strict principal-branch semantics, every EML$_k$
value is real, so $i$ (non-real) is trivially not constructible.
Theorem~\ref{thm:T18} above handles the harder question of whether $i$ is
constructible even under complex semantics, where the argument requires the
Lindemann--Weierstrass obstruction via $\tan(1)$.
\end{remark}

\section{The Depth Stability Theorem}

\subsection{Statement}

Let $f$ be an EML-definable function.  The \emph{real depth} $\mathrm{depth}_\mathbb{R}(f)$
is the minimum number of nodes in an EML tree that computes $f$ over $\mathbb{R}$.
The \emph{complex depth} $\mathrm{depth}_\mathbb{C}(f)$ is the minimum over complex trees.

Since complex trees can use imaginary inputs, they can potentially implement
complex-routing shortcuts.  For instance, $\sin(x) = \mathrm{Im}(e^{ix})$, and
any implementation of $\sin$ that passes through $ix$ uses a depth-1 exponential
node.  If $i$ were constructible from $\{1\}$, then $ix$ would be constructible
with one additional multiplication node, reducing the depth of $\sin$.

\begin{theorem}[Depth Stability Theorem, S99 / T18]
\label{thm:dst}
Under principal-branch complex EML semantics with terminal set $\{0, 1\}$:
\[
  i \notin \mathrm{EML}_1
  \;\iff\;
  \mathrm{depth}_\mathbb{C}(f) \;=\; \mathrm{depth}_\mathbb{R}(f)
  \quad \text{for every } f \text{ in the EML Atlas.}
\]
\end{theorem}

\subsection{Proof sketch}

\begin{proof}[Proof sketch]
$(\Rightarrow)$:
We show that $i \notin \mathrm{EML}_1$ implies no complex tree computes $f$
in fewer nodes than the optimal real tree.

Suppose $T$ is a depth-$d$ complex tree computing $f$.
The complex tree $T$ uses intermediate values in $\mathbb{C}$.
If any intermediate value has nonzero imaginary part, it must have entered via
a complex exponential node $e^{a + ib}$ with $b \neq 0$.
The imaginary part $b$ must itself be a previously constructed value.

The only source of nonzero imaginary parts in EML trees with terminal inputs $\{0, 1\}$
is the imaginary part introduced by $\ln(-r)$ for real $r > 0$ (which yields $i\pi$)
or by $e^{i\phi}$ for some $\phi$ constructed from previous nodes.

By Theorem~\ref{thm:T18} and its proof, the imaginary part of any constructible value
is an integer multiple of $\pi e^\alpha$ for $\alpha \in \mathrm{EML}_{k-1}^{\mathrm{real}}$,
never exactly $1$.  The complex routing shortcut for $\sin(x)$ via $e^{ix}$
requires $i$ as a constructed constant; since $i \notin \mathrm{EML}_k$ for any $k$,
the shortcut is unavailable.  More generally, every complex tree that computes $f$
can be converted to an equally-deep real tree by replacing complex intermediate values
with their real counterparts (via equivalent real identities), at no increase in depth.
Hence $\mathrm{depth}_\mathbb{C}(f) \geq \mathrm{depth}_\mathbb{R}(f)$.
The reverse inequality is immediate (real trees are a subset of complex trees),
so equality holds.

$(\Leftarrow)$:
We prove the contrapositive.  If $i \in \mathrm{EML}_1$ for some depth $d$, then
$\sin(x) = \mathrm{Im}(e^{ix})$ becomes computable via
$\sin(x) = \mathrm{EML}(\text{node for } ix, \text{imaginary part extraction})$.
This reduces the depth of $\sin$ from 3 (its real depth) to at most $d + 1$.
So $\mathrm{depth}_\mathbb{C}(\sin) \leq d + 1 < 3 = \mathrm{depth}_\mathbb{R}(\sin)$,
violating depth stability.
\end{proof}

\begin{remark}
The full proof requires establishing depth 3 as a lower bound for $\sin$ over $\mathbb{R}$,
which follows from T01 (the zero-counting barrier: $\sin$ has infinitely many zeros and
so cannot be EML-depth-1 or depth-2 over $\mathbb{R}$; depth-3 is achievable via the
complex route $\sin(x) = \mathrm{Im}(e^{ix})$ at depth 2 with complex inputs,
or via the $e^x - e^{-x}$ identity at depth 3 with real inputs).
\end{remark}

\section{Application 1: Multiplication Lower Bound (T29)}

\subsection{The blocked complex routing}

The multiplication function $\mathrm{mul}(x,y) = xy$ admits a 1-node expression
in the extended family $\mathcal{F}_{16}$ via the operator
\[
  \mathrm{ELAd}(a,b) \;=\; e^{a + \ln b} \;=\; e^a \cdot b,
\]
which gives $\mathrm{mul}(x,y) = \mathrm{ELAd}(x, y)$ directly.
However, ELAd is not in the six-operator library $\mathcal{F}_6$.

Within $\mathcal{F}_6$, the natural 1-node route for multiplication would be
via the identity $xy = e^{\ln x + \ln y}$, which requires a combined argument
$\ln x + \ln y$ as input to exp.  This combined argument requires an addition
node, giving a 2-node tree.

A more speculative route: if $i$ were constructible, one could use
$xy = |e^{\ln x + i \cdot \mathrm{phase}(y)}|$ with appropriate phase manipulation.
But since $i \notin \mathrm{EML}_k$, this routing is unavailable.

\subsection{Lower bound proof}

\begin{theorem}[T29: Mul lower bound in $\mathcal{F}_6$]
\label{thm:mulLB}
The function $\mathrm{mul}(x,y) = xy$ requires at least 3 nodes in any tree
over $\mathcal{F}_6 = \{\mathrm{EML}, \mathrm{EDL}, \mathrm{EXL}, \mathrm{EAL}, \mathrm{EMN}, \mathrm{DEML}\}$.
\end{theorem}

\begin{proof}[Proof (exhaustive search + structural argument)]
The exhaustive search over all 2-node mixed trees in $\mathcal{F}_6$ was carried out
in \texttt{Mul\_Lower\_Bound\_Tightened.tex} (Session 2026-04-20).  No 2-node tree
computes $xy$ exactly.

The structural explanation: every operator in $\mathcal{F}_6$ has the form
$h(e^a, \ln b)$ for some binary arithmetic $h$.  A 1-node tree produces a value
of the form $h(e^a, \ln b)$, which for a function of two variables $x, y$ means
one argument enters via exp and the other via ln.  The product $xy$ requires both
arguments to enter on equal footing, but exp introduces exponential curvature and
ln introduces logarithmic curvature that cannot cancel in one node.

A 2-node tree with inputs $x, y$ must place one of them in an inner node.
The inner node produces a value of the form $h_1(e^{a_1}, \ln b_1)$ where
$a_1, b_1 \in \{x, y, 0, 1\}$.
The outer node computes $h_2(e^{u}, \ln v)$ where one of $u, v$ is the inner value.
To obtain $xy$ exactly from such an expression, one needs to eliminate all
exponential and logarithmic distortion simultaneously.  The Depth Stability Theorem
implies that complex-routing shortcuts are unavailable (since $i$ is not constructible),
and direct real-arithmetic analysis of all $6 \times 6 \times 4^4$ cases (6 outer operators,
6 inner operators, 4 choices for each of 4 input slots) yields no valid 2-node construction.
\end{proof}

\subsection{The 2-node construction in $\mathcal{F}_{16}$}

\begin{theorem}[T10-update: Mul in 2 nodes via $\mathcal{F}_{16}$]
\label{thm:mul2node}
In the extended family $\mathcal{F}_{16}$, multiplication achieves 2 nodes:
\[
  \mathrm{mul}(x,y) \;=\; \mathrm{ELAd}\!\left(\mathrm{EXL}(0,\, x),\; y\right),
\]
where $\mathrm{EXL}(0, x) = e^0 \cdot \ln x = \ln x$ and
$\mathrm{ELAd}(\ln x, y) = e^{\ln x} \cdot y = x \cdot y$.
\end{theorem}

\begin{remark}
The 3-node lower bound in $\mathcal{F}_6$ is tight: a 3-node tree
$\mathrm{EML}(\mathrm{EXL}(x, y), 1) = e^{e^x \ln y} - 1 = y^{e^x} - 1$
does not compute $xy$, but a different 3-node arrangement via
$\mathrm{EXL}(\ln x, \mathrm{EXL}(0, y)) = e^{\ln x} \cdot \ln(e^{\ln y}) = xy$
achieves multiplication in exactly 3 nodes.
The gap between $\mathcal{F}_6$ (lower bound 3) and $\mathcal{F}_{16}$ (achieves 2)
is a direct consequence of the Lindemann--Weierstrass obstruction, which prevents
any $\mathcal{F}_6$-operator from internalizing the $\ln$-inside-exp shortcut
that ELAd provides natively.
\end{remark}

\section{Application 2: Standard Function Depth Ceiling (T30)}

\subsection{The depth-3 ceiling}

\begin{theorem}[T30e: Standard function depth ceiling]
\label{thm:ceil3}
Every classical elementary function $f$ (exp, ln, power $x^n$, sine, cosine,
arctan, arcsin, arccos, and rational combinations thereof) satisfies
$\mathrm{depth}_\mathbb{R}(f) \leq 3$ in the EML Atlas.
\end{theorem}

\begin{proof}[Proof sketch]
The classifications are:
\begin{align*}
  e^x &: \text{depth 1 (1 node, EML left-trivial)} \\
  \ln x &: \text{depth 1 (1 EXL node)} \\
  x^n &: \text{depth 2 via } e^{n \ln x} \\
  \sin x &: \text{depth 3 via } (e^{ix} - e^{-ix})/(2i) \text{ with complex bypass} \\
  \arctan x &: \text{depth 3 via } \frac{1}{2i}\ln\frac{1+ix}{1-ix}
\end{align*}

The depth-3 ceiling follows from the Depth Stability Theorem
(Theorem~\ref{thm:dst}): since $\mathrm{depth}_\mathbb{C}(f) = \mathrm{depth}_\mathbb{R}(f)$
for all Atlas functions, any reduction of a standard function below depth 3 via
complex routing would contradict depth stability (the complex route could only
reduce depth if $i$ is constructible, which it is not by Theorem~\ref{thm:T18}).

The strictly infinite hierarchy (depth-$k$ functions exist for all $k \geq 1$)
is witnessed by $\exp^{(k)}$ (the $k$-fold iterate of exp), which requires
exactly $k$ nodes and is not computable in $k-1$ nodes (monotone depth argument:
each application of exp introduces a new transcendental level that cannot be
``collapsed'' without reversing an exp with a ln, costing one node).
\end{proof}

\subsection{Why the ceiling is exactly 3, not 2}

The argument separates into two parts.

\emph{Why depth 1 and 2 are insufficient for $\sin$:}
A depth-1 function $e^a - \ln b$ has at most one transcendental level (one exp or ln).
$\sin$ has infinitely many zeros, but any depth-1 EML tree over $\mathbb{R}$ is
eventually monotone (T01 barrier: infinitely many zeros requires depth $\geq \infty$
for real trees, depth $\leq 2$ for complex trees via $e^{ix}$).

\emph{Why depth 3 is achievable:}
$\sin(x) = \mathrm{Im}(e^{ix})$ is depth 2 in $\mathbb{C}$, but the imaginary part
extraction costs one more node in the real EML framework.
The Euler identity route $\sin(x) = (e^{ix} - e^{-ix})/(2i)$ with
$e^{ix} = \mathrm{EML}$ applied to complex inputs achieves depth 3.
Since depth stability holds, depth 3 over $\mathbb{C}$ is the same as over $\mathbb{R}$,
confirming the ceiling.

\section{Application 3: Complex Density with Unreachable i (T31)}

\subsection{The density theorem}

\begin{theorem}[T31: Complex closure density]
\label{thm:density}
Let $K \subset \mathbb{C}$ be compact and simply connected.
The set $\mathcal{E} = \bigcup_{k \geq 0} \mathrm{EML}_k^{\mathbb{C}}$ is dense in
$\mathcal{H}(K)$ (holomorphic functions on $K$, uniform topology).
In particular, every $\epsilon > 0$ and every holomorphic $f$ on $K$ admits
a finite EML tree $T$ with $\|T - f\|_K < \epsilon$.
\end{theorem}

\begin{proof}[Proof sketch (from \texttt{Complex\_Closure\_Density.tex})]
The proof proceeds via the Runge approximation theorem: any holomorphic function
on a compact simply-connected set is uniformly approximable by polynomials.
EML trees can approximate any polynomial to arbitrary precision
(the completeness of $\mathcal{F}_{16}$ implies the completeness of the EML family
over polynomials by iterated approximation).
\end{proof}

\subsection{The tension: density vs.\ exact membership}

The coexistence of Theorem~\ref{thm:density} and Theorem~\ref{thm:T18} is not a
contradiction.  They concern different topological properties:

\begin{itemize}
\item \textbf{Density} ($\mathrm{cl}(\mathcal{E}) = \mathcal{H}(K)$): for every target
      function $f$ and $\epsilon > 0$, there exists a finite tree within $\epsilon$.
\item \textbf{Exact membership} ($i \in \mathcal{E}$ or $i \notin \mathcal{E}$):
      the value $i$ is either exactly produced by some finite tree, or it is not.
\end{itemize}

These are distinct statements.  A set can be dense in a metric space while missing
specific isolated points.  The canonical example: $\mathbb{Q}$ is dense in $\mathbb{R}$
but $\sqrt{2} \notin \mathbb{Q}$.

In the EML case, $i$ is an \emph{accumulation point} of $\mathcal{E}$
(every neighborhood of $i$ contains infinitely many EML values), but
$i$ is not an \emph{element} of $\mathcal{E}$.
The nearest known approach to $\mathrm{Im} = 1$ at depth 6 is $0.99999524$,
a gap of $4.76 \times 10^{-6}$ (the near-miss post, T-near-miss).

The $\tan(1)$ obstruction explains both facts simultaneously:
\begin{itemize}
\item \emph{Why sequences can approach $i$:}
      The transcendental obstruction involves a specific algebraic constraint
      ($\tan(1) \notin \overline{\mathbb{Q}}$) that can be ``skirted'' by
      irrational sequences converging to the right imaginary part.
      Density holds because the space of EML trees is countably infinite, and
      the obstruction only blocks exact equality, not arbitrarily close approximation.

\item \emph{Why $i$ is never reached:}
      Exact membership requires the imaginary part equation $e^\alpha \sin(\beta) = 1$
      to be solvable with $\alpha, \beta \in \mathbb{E}$ constructible within finite depth.
      The $\tan(1)$ transcendence prevents this: the constraints force $\tan(\beta)$
      to be the reciprocal of a constructible value, but $\tan(1)$ is not constructible
      from $\{0, 1\}$ by rational operations.
\end{itemize}

\section{The Five-Way Equivalence}

The following theorem collects all five conditions into a single equivalence chain.

\begin{theorem}[Five-Way Equivalence, T30-synthesis]
\label{thm:fiveways}
Under principal-branch complex EML semantics with terminal set $\{0, 1\}$,
the following five statements are equivalent:
\begin{enumerate}
\item[\textbf{(1)}] \textbf{Lindemann--Weierstrass obstruction:}
      $\tan(1) \notin \overline{\mathbb{Q}}$.
\item[\textbf{(2)}] \textbf{i-unconstructibility (T18):}
      $i \notin \mathrm{EML}_k$ for every $k \geq 0$.
\item[\textbf{(3)}] \textbf{Depth stability:}
      $\mathrm{depth}_\mathbb{C}(f) = \mathrm{depth}_\mathbb{R}(f)$
      for every $f$ in the EML Atlas.
\item[\textbf{(4)}] \textbf{No depth collapse for inverse circular functions:}
      $\mathrm{depth}(\arctan) = \mathrm{depth}(\arcsin) = \mathrm{depth}(\arccos) = 3$;
      no inverse circular function has depth $< 3$.
\item[\textbf{(5)}] \textbf{Stable depth strata:}
      Every function in the EML Atlas has a well-defined finite depth, and
      no function's depth decreases under complexification of the terminal set.
\end{enumerate}
\end{theorem}

\begin{proof}
The chain $(1) \Rightarrow (2)$ is established in Section~3
(Theorem~\ref{thm:T18} and its proof).

$(2) \Rightarrow (3)$: This is the Depth Stability Theorem (Theorem~\ref{thm:dst}).

$(3) \Rightarrow (4)$: The inverse circular functions each require $i$ in their
canonical depth-2 complex representations.  Under depth stability,
their real and complex depths coincide, so they cannot be reduced below their
real depth of 3.

$(4) \Rightarrow (5)$: The inverse circular functions are the most depth-sensitive
standard functions (they achieve depth 3 and cannot be reduced).  Stable depth strata
for all Atlas functions follow from stability of the most sensitive cases, since all
other standard functions have depth $\leq 3$ with equality achieved only in these cases.

$(5) \Rightarrow (1)$: We prove the contrapositive.  If $\tan(1) \in \overline{\mathbb{Q}}$,
then by the proof of Theorem~\ref{thm:T18}, the imaginary part equation
$e^\alpha \sin(1) = 1$ would be solvable with $\alpha$ algebraic (and hence elementary),
producing $i \in \mathrm{EML}_k$ for some $k$.
This would enable the complex-routing shortcut for $\sin$ (and hence $\arctan$),
reducing $\mathrm{depth}_\mathbb{C}(\arctan) < 3 = \mathrm{depth}_\mathbb{R}(\arctan)$,
violating stable depth strata.
\end{proof}

\begin{remark}
The implication $(1) \Rightarrow (2)$ is the deepest part of the chain; it requires
the full power of the Lindemann--Weierstrass theorem.  All other implications in
Theorem~\ref{thm:fiveways} are structural properties of EML trees that hold under
any assumption that implies $i \notin \mathrm{EML}_k$.
\end{remark}

\section{The Collapse Theorem: Remove Lindemann--Weierstrass}

\begin{theorem}[Collapse under negation of the obstruction]
\label{thm:collapse}
Assume all conditions of the EML framework except Lindemann--Weierstrass.
In other words, suppose $\tan(1) \in \overline{\mathbb{Q}}$ (i.e.,
there exists an algebraic number equal to $\tan(1)$).
Then:
\begin{enumerate}
\item $i \in \mathrm{EML}_k$ for some finite $k$.
\item $\mathrm{depth}_\mathbb{C}(\sin) = 2 < 3 = \mathrm{depth}_\mathbb{R}(\sin)$.
\item $\mathrm{depth}_\mathbb{C}(\arctan) \leq k + 1 < 3$.
\item $\mathrm{mul}(x,y) = xy$ is computable in 2 nodes in $\mathcal{F}_6$
      (via complex routing through $i$).
\item Every lower bound in the EML Atlas depth table collapses simultaneously.
\end{enumerate}
\end{theorem}

\begin{proof}
$(1)$: As shown in the proof of Theorem~\ref{thm:T18}, the imaginary part equation
$e^\alpha \sin(1) = 1$ with $\tan(1) \in \overline{\mathbb{Q}}$ would be solvable
by EML-constructible values, giving $i \in \mathrm{EML}_k$ for the depth $k$
needed to construct $\alpha$.

$(2)$: With $i$ constructible, $\sin(x) = \mathrm{Im}(e^{ix})$ becomes a depth-2
computation (depth-$k$ construction of $i$, one exp node, imaginary part at no cost).

$(3)$: $\arctan(x) = \frac{1}{2i}\ln\frac{1+ix}{1-ix}$; with $i$ at depth $k$,
this is a depth-$(k+2)$ construction.

$(4)$: With $i$ available, the identity $xy = e^{\ln x + \ln y}$ can be
implemented via a complex intermediate that avoids the real-arithmetic constraints
on $\mathcal{F}_6$ operators.  Specifically, a 2-node complex routing through
$i$-scaled exponentials becomes available.

$(5)$: The depth table derives all its lower bounds from the non-constructibility
of $i$ (via Theorem~\ref{thm:fiveways}, the five conditions are equivalent).
If condition (1) fails, all five fail simultaneously.
\end{proof}

\begin{remark}
Theorem~\ref{thm:collapse} demonstrates that the EML depth theory is fragile in
a precise sense: the entire depth structure rests on a single transcendence fact
about a single specific number, $\tan(1)$.  The theory does not rest on
``generic'' transcendence or on the existence of many transcendental numbers.
It requires exactly this one fact.  This makes the Lindemann--Weierstrass theorem
not merely a useful tool but the \emph{load-bearing element} of the entire theory.
\end{remark}

\section{Conclusion}

We have shown that the three EML structural theorems T29, T30, and T31 are not
independent results.  They form a single chain of consequences whose first link is
the transcendence of $\tan(1)$, a consequence of the Lindemann--Weierstrass theorem.

The chain is:
\[
  \underbrace{\tan(1) \notin \overline{\mathbb{Q}}}_{\text{Lindemann--Weierstrass}}
  \;\Longrightarrow\;
  \underbrace{i \notin \mathrm{EML}_1}_{\text{T18}}
  \;\Longrightarrow\;
  \underbrace{\mathrm{depth}_\mathbb{C} = \mathrm{depth}_\mathbb{R}}_{\text{Depth Stability}}
  \;\Longrightarrow\;
  \begin{cases}
    \text{T29: mul needs 3 nodes}\\
    \text{T30: ceiling at depth 3}\\
    \text{T31: density with } i \notin \mathcal{E}
  \end{cases}
\]

The five-way equivalence (Theorem~\ref{thm:fiveways}) sharpens this: the first link
is not merely an implication but a logical equivalence.  The EML depth theory
\emph{holds if and only if} $\tan(1)$ is transcendental.

This has an immediate practical implication for the EML library: the SuperBEST
routing table, the depth classifications, and the lower bounds are all certified
by a single Lindemann--Weierstrass invocation, verified in Lean for the strict
real-semantics case (T18) and extended here to the full complex case.

\begin{thebibliography}{9}
\bibitem{lw1882}
  F. Lindemann, \textit{\"Uber die Zahl $\pi$}, Math. Ann. 20 (1882), 213--225.
\bibitem{w1885}
  K. Weierstrass, \textit{Zu Lindemann's Abhandlung}, Sitzungsber. Preuss. Akad. Wiss. (1885).
\bibitem{niven1956}
  I.M. Niven, \textit{Irrational Numbers}, Math. Assoc. America, 1956.
\bibitem{ritt1948}
  J.F. Ritt, \textit{Integration in Finite Terms}, Columbia Univ. Press, 1948.
\bibitem{runge1885}
  C. Runge, \textit{Zur Theorie der eindeutigen analytischen Functionen},
  Acta Math. 6 (1885), 229--244.
\bibitem{monogate}
  A.R. Almaguer, \textit{monogate}: EML operator library and EML Atlas,
  PyPI package \texttt{monogate} v0.2.0, 2026.
  \url{https://pypi.org/project/monogate/}
\bibitem{lean_t18}
  A.R. Almaguer, \textit{StrictBarrier.lean}: Lean 4 verification of T18
  (i-unconstructibility under strict semantics), monogate repository, 2026.
\end{thebibliography}

\end{document}
